\section{A three-node spectral primer}
\label{sec:book-spectral-primer}

Before applying spectral transforms to protein operators, consider the path on three positions,
\begin{equation}
A=
\begin{pmatrix}
0&1&0\\
1&0&1\\
0&1&0
\end{pmatrix},
\qquad
D=\operatorname{diag}(1,2,1),
\qquad
L=D-A=
\begin{pmatrix}
1&-1&0\\
-1&2&-1\\
0&-1&1
\end{pmatrix}.
\label{eq:book-three-node-laplacian}
\end{equation}
For a scalar signal $u=(u_1,u_2,u_3)^{\mathsf T}$ on the positions,
\begin{equation}
u^{\mathsf T}Lu=(u_1-u_2)^2+(u_2-u_3)^2.
\label{eq:book-three-node-dirichlet}
\end{equation}
The quadratic form is small when neighboring positions carry similar values.  The eigenvalues are
$0,1,3$.  The zero mode is constant, the middle mode contrasts the two endpoints, and the highest
mode alternates across both edges.  Low eigenvalue therefore means smoothness relative to this
declared graph and this declared node signal.

Three choices precede every biological interpretation of these modes.  First, the adjacency must
be supplied: contact, covariance, attention, and an operator norm produce different $A$.  Second,
the node signal must be typed: a scalar position score, an amino-acid contrast, and an atomic
displacement live in different spaces.  Third, any null subtraction, symmetrization, or
normalization changes the operator being diagonalized.

This primer uses the unnormalized Laplacian.  Later sections introduce normalized variants; they
rescale node coordinates while preserving the same edgewise smoothness principle.  For
nonsymmetric routes, eigenvectors alone are insufficient, so stationary geometry and left--right
structure enter the audit.

The rest of the chapter generalizes this calculation in stages.  Directed routes require a
stationary geometry; signed edges retain a declared polarity; connection Laplacians transport
categorical vectors between sites; wavelets localize selected spectral bands.  The three-node path
is the baseline against which each additional structure should be read.
